\chapter{Conclusion and Outlook}\label{Chap:Conclusion}The primary objective of this thesis was to develop and implement a robust framework for evaluating data consistency in reaction model optimization. To this end, a reaction rate coefficient uncertainty package was developed to estimate model uncertainties, and an ignition delay time uncertainty package was implemented to quantify experimental data uncertainties. These components were integrated into the Bound-to-Bound Data Collaboration (B2BDC) framework to perform comprehensive consistency analyses. To support those functionalities, Cantera ignition delay time simulations were implemented. Sensitivity analysis was implemented to determine active reactions and reduce computational cost of the rest of the analysis. The results of these analyses are heavily dependent on the quality and availability of the underlying data. Although literature sources for ignition delay time often report critical parameters such as $T_5$, $P_5$, $\phi$ uncertainties, or pressure rise, this information is rarely integrated into large experimental databases. A similar challenge exists regarding kinetics databases. While the NIST Chemical Kinetics Database \cite{j_a_manion_r_e_huie_r_d_levin_d_r_burgess_jr_v_l_orkin_w_tsang_w_s_mcgivern_j_w_hudgens_v_d_knyazev_d_b_atkinson_e_chai_a_m_tereza_c-y_lin_t_c_allison_w_g_mallard_f_westley_j_t_herron_r_f_hampson_and_d_h_frizzell_nist_nodate} is highly comprehensive, it cannot be accessed programmatically. Conversely, the RMG database \cite{johnson_rmg_2022-1} is accessible and user-friendly but often lacks essential information, such as specific temperature ranges for kinetic data.Future work should focus on the following areas:\begin{itemize}\item Full integration of the packages developed in this thesis into the TUMKIN infrastructure to streamline the optimization workflow.\item Building a comprehensive database for TUMKIN that supports the analysis and optimization of a wider range of reaction models and fuels.\item Extending the uncertainty quantification methods used for ignition delay times to other experimental targets, such as laminar flame speeds and species concentrations.\item Exploring the implementation of model discrepancy functions to explicitly account for inherent model inadequacies both during consistency analysis and the optimization process.\end{itemize}







